\chapter*{Résumé}
\addcontentsline{toc}{chapter}{Résumé}

\begin{onehalfspace}
Ce projet présente la conception et la mise en œuvre d'une solution de Business Intelligence (BI) dédiée à l'analyse de l'attrition et de la performance des employés. L'objectif principal est d'identifier les facteurs déterminants qui influencent le départ des collaborateurs au sein d'une organisation afin de proposer des stratégies de rétention basées sur les données. En s'appuyant sur un jeu de données de ressources humaines comprenant 1470 observations, nous avons déployé une architecture décisionnelle complète. La méthodologie adoptée comprend une phase d'exploration des données (EDA), la conception d'un modèle dimensionnel en étoile optimisé pour les analyses RH, ainsi qu'un processus ETL rigoureux pour le nettoyage et la transformation des informations. Les résultats, présentés sous forme de tableaux de bord interactifs, mettent en évidence l'impact prépondérant des heures supplémentaires, du niveau de rémunération et de l'équilibre vie-travail sur la stabilité des effectifs. Ce travail offre une base analytique solide pour orienter les politiques de gestion des talents et ouvre la voie à l'utilisation de modèles prédictifs pour anticiper les risques d'attrition.

\textbf{Mots-clés :} Business Intelligence, HR Analytics, Attrition, Modélisation dimensionnelle, ETL, Visualisation de données.
\end{onehalfspace}
